% Lesson: the limit of a function at a point (English).

\begin{frame}
\didactatitle{The limit of a function at a point}

\bymedium{%
  $f(x)$ gets close to $L$ as $x$ gets close to $a$. The definition says what
  ``gets close'' means.
}{%
  The intuitive idea is that $f(x)$ gets as close as we like to $L$ as long as
  $x$ is close enough to $a$. The definition turns that sentence into an
  inequality that can be checked.
}

\begin{definition}[Limit]
Let $f$ be a function defined on an open interval containing $a$, except
possibly at $a$ itself. We say that $f$ has \keyterm{limit} $L$ at $a$, and
write $\lim_{x \to a} f(x) = L$, if for every $\varepsilon > 0$ there is a
$\delta > 0$ such that
\begin{keyformula}
  0 < |x - a| < \delta \implies |f(x) - L| < \varepsilon .
\end{keyformula}
\end{definition}

\begin{notesonly}
The condition $0 < |x - a|$ leaves out the point $a$ itself: the limit does
not depend on the value of $f$ at $a$, nor even on $f$ being defined there. It
only looks at what happens \emph{near} $a$. That is why it makes sense to ask
for the limit of $(x^2 - 1)/(x - 1)$ at $x = 1$, where the quotient is not
defined.

Moreover, the limit, if it exists, is unique: two different values
$L \neq L'$ cannot both be within $\varepsilon = |L - L'|/2$ of $f(x)$.
\end{notesonly}

\begin{teaching}[Pacing]
Twenty minutes. Draw the band $(L - \varepsilon, L + \varepsilon)$ before
writing the formula: the order of the quantifiers (first you are given
$\varepsilon$, then you choose $\delta$) is easier to see on the board.
\end{teaching}
\end{frame}
